A análise de sobrevivência constitui um dos principais ramos da estatística aplicada voltado ao estudo do
tempo até a ocorrência de um evento de interesse, como falha de um sistema, recorrência de uma doença ou
morte de um indivíduo. Uma característica central desse tipo de dado é a presença de censura, isto é,
situações nas quais o tempo exato do evento não é completamente observado. Esse aspecto impõe desafios
metodológicos relevantes, exigindo o desenvolvimento de modelos capazes de lidar adequadamente com
informações incompletas.

Além da presença de censura, em diversos contextos práticos observa-se a existência de uma fração de
indivíduos que não experimentará o evento de interesse, mesmo sob longos períodos de acompanhamento.
Essa característica motivou o desenvolvimento dos chamados modelos de fração de cura, que buscam representar
explicitamente essa parcela da população.

Tradicionalmente, modelos como o de Cox têm sido amplamente utilizados nesse contexto, devido à sua
flexibilidade semiparamétrica e interpretabilidade. No entanto, tais abordagens assumem, em geral, uma
estrutura linear no preditor do risco, o que pode ser inadequado em cenários nos quais a relação entre as
covariáveis e o tempo de sobrevivência apresenta padrões não lineares ou de maior complexidade.

Nesse cenário, métodos baseados em kernels surgem como alternativas promissoras, pois permitem incorporar
flexibilidade ao processo de modelagem por meio de mapeamentos implícitos para espaços de maior dimensão.
Essa abordagem possibilita capturar estruturas não lineares sem a necessidade de especificar explicitamente
a forma funcional da relação entre covariáveis e sobrevivência.

O presente trabalho propõe um novo modelo semiparamétrico para dados de sobrevivência com fração de cura,
denominado GG-KM, que combina a flexibilidade de métodos baseados em kernel com a estrutura do modelo de
tempo de promoção. O objetivo principal é investigar o desempenho do modelo sob diferentes especificações
de kernel, avaliando sua capacidade preditiva e sua adequação em cenários com diferentes graus de
complexidade.

Para isso, são conduzidos experimentos em dados simulados, nos quais é possível controlar a estrutura geradora,
e em bases de dados reais amplamente utilizadas na literatura. A avaliação é realizada principalmente por
meio do Integrated Brier Score (IBS), permitindo analisar a qualidade preditiva do modelo e a influência
da escolha do kernel sob diferentes cenários.

\section{Organização do Trabalho}

Este trabalho está estruturado da seguinte forma. No Capítulo 4, são apresentados os conceitos fundamentais
de análise de sobrevivência, incluindo noções de censura, funções de sobrevivência, modelo de Cox, além de descrever brevemente modelos com fração de cura.
No Capítulo 3, são introduzidos alguns conceitos que circundam métodos baseados em kernel. Dessa forma, são apresentadas
definições como espaços de Hilbert, kernels positivos definidos e alguns kernels principais.

No Capítulo 5, é detalhada a metodologia adotada, abrangendo o particionamento dos dados, a geração de dados
simulados com fração de cura. O Capítulo 6 apresenta e discute os resultados obtidos tanto em dados simulados quanto em dados
reais. Por fim, o Capítulo 7 apresenta as conclusões do trabalho.
